\section*{Abstract}
\addcontentsline{toc}{section}{Abstract}

Occupational choice is commonly made with limited direct knowledge of the work involved. Research on career decision-making identifies insufficient occupational information as a principal source of difficulty, and the realistic job preview literature establishes that balanced, candid previews improve person--job fit and reduce subsequent regret. Delivering such previews has historically required either workplace exposure or content authored occupation by occupation, neither of which scales. This project designs, implements, and evaluates a platform that uses generative language models to deliver experiential occupational exposure at greater breadth, and establishes whether competence demonstrated within such an experience can be measured with sufficient reliability to be useful.

The report establishes the evidence base for this approach, reviews existing research prototypes and deployed systems against criteria derived from properties the literature identifies as consequential for learning systems, and specifies a proposed platform in which occupations are presented as performable tasks, responses are assessed against expert-authored rubrics, and demonstrated competence determines subsequent exploration. The review finds that the constituent techniques have been demonstrated separately but not in combination, and that existing systems in this area report assessment agreement without chance correction or difficulty stratification. The proposed evaluation addresses this through an expert-reviewed reference set constructed within a single occupational field. The report covers Chapters 1 to 5.1: introduction, background, related work, the proposed system, and its analysis.

\textbf{Keywords:} career exploration, experiential learning, realistic job preview, large language models, automated assessment, stealth assessment, serious games.
